\subsubsection{Sistema Operativo}

La decisión de utilizar FreeRTOS como sistema operativo para el desarrollo de este proyecto se fundamenta en varias consideraciones claves. Primero, FreeRTOS está integrado de manera nativa en el framework ESP-IDF de Espressif, lo que facilita enormemente el desarrollo y la integración de funcionalidades en tiempo real. Al utilizar la API de Espressif, se aprovechan componentes y servicios que ya están optimizados para trabajar con FreeRTOS, lo que se traduce en una mayor estabilidad y rendimiento en aplicaciones críticas.

FreeRTOS permite una gestión eficiente de tareas, permitiendo la concurrencia y la sincronización de procesos de manera robusta y controlada. Esto es esencial para aplicaciones que requieren una respuesta rápida a eventos externos, como la gestión de sensores, la comunicación inalámbrica y la realización de múltiples operaciones simultáneamente. La modularidad y escalabilidad que ofrece FreeRTOS permite desarrollar soluciones complejas sin incurrir en sobre costos de recursos, lo que es especialmente valioso en sistemas embebidas con limitaciones de hardware. \cite{barry2016mastering}

El uso de FreeRTOS está respaldado por un amplio ecosistema de documentación y soporte, facilitando el desarrollo, la depuración y el mantenimiento de aplicaciones en el microcontrolador ESP32. Esta sinergia entre el hardware y el sistema operativo garantiza una integración fluida y optimizada para las exigencias de aplicaciones modernas en sistemas embebidos y soluciones IoT, lo que se ajusta perfectamente al desarrollo de este proyecto.

\begin{figure}[H]
  \centering
  \includegraphics[width=0.4\linewidth]{images/freertos.jpg}
  \caption{Sistema operativo usado por los microcontroladores}
  \label{fig:sistema_operativo}
\end{figure}